\documentclass[a4paper,12pt,titlepage,final]{article}

\usepackage[T1,T2A]{fontenc}     % форматы шрифтов
\usepackage[utf8]{inputenc}      % кодировка символов, используемая в этом файле
\usepackage[russian]{babel}      % пакет русификации
\usepackage{tikz}                % для создания иллюстраций
\usepackage{geometry}            % для настройки размера полей
\usepackage{indentfirst}         % для отступа в первом абзаце секции
\usepackage{booktabs}
\usepackage{ltablex}
\usepackage{float}
\usepackage{blindtext}

\geometry{a4paper,left=30mm,top=30mm,bottom=30mm,right=30mm}

\renewcommand{\thesection}{\arabic{section}.}

\newcommand{\university}{Московский государственный университет \\имени М.~В.~Ломоносова\\
	Факультет вычислительной математики и кибернетики}

\begin{document}

\begin{titlepage}
    \begin{center}
	{\large \sc \university\\}
	\vfill
	{\Large \sc Отчет по заданию №1}\\
	~\\
	{\large \bf <<Методы сортировки>>}\\ 
	~\\
	{\large \bf Вариант 3 / 3 / 1 + 4}
    \end{center}
    \begin{flushright}
	\vfill {Выполнил:\\
	студент 103 группы\\
	Огай~В.~Л.\\
	~\\
	Преподаватель:\\
	Кузьменкова~Е.~А.}
    \end{flushright}
    \begin{center}
	\vfill
	{\small Москва\\2024}
    \end{center}
\end{titlepage}

\tableofcontents

\newpage
\section{Постановка задачи}
Требуется реализовать два метода сортировки массивов: метод «пузырька» и быструю сортировку (рекурсивная реализация) - чисел типа double в порядке неубывания модулей и провести их экспериментальное сравнение. В экспериментальное сравнение помимо реализации вышеперечисленных методов входит подсчёт произведённых сравнений и перестановок в массивах в ходе сортировок. Эксперимент проводится на четырех массивах: отсортированный в прямом порядке (1), отсортированный в обратном порядке (2), случайно организованный 1 (3), случайно организованный 2 (4).

\newpage
\section{Результаты экспериментов}
\textbf{\vspace{0.2cm} \hspace{-0.92cm} \large Теоретические оценки\\}
\vspace{0.2cm}
Пусть $N$ -- количество элементов сортируемого массива.\\
\vspace{0.2cm}
\textbf{Сортировка методом «пузырька»} \\
\vspace{0.15cm}
\hspace{-0.32cm} Число сравнений: $\frac{1}{2}N(N - 1)$\\
Среднее количество перемещений в данном методе пропорционально $\frac{1}{4}(N^2 - N)$.(\cite{first}, разд. 5.2.2)

\hspace{-0.6cm} \textbf{\vspace{0.2cm} Быстрая сортировка, рекурсивная реализация\\}
\hspace{-0.32cm} Число сравнений: $N\log{N}$\\
Среднее количество перемещений в данном методе пропорционально $\frac{N}{6}\log{N}$.(\cite{second}, разд. 2.2.6)

\textbf{\hspace{-0.8cm} \vspace{0.2cm} \large Экспериментальные значения\\}
В таблицах 1-4 номера массивов соответствуют массивам со следующими видами расстановки элементов:
\begin{itemize}
    \item Массив 1: массив, отсортированный в прямом порядке
    \item Массив 2: массив, отсортированный в обратном порядке
    \item Массив 3: массив, элементы которого идут в случайном порядке 1
    \item Массив 4: массив, элементы которого идут в случайном порядке 2
\end{itemize}

\begin{table}[H]
\centering
\begin{tabular}{*{6}{>{\raggedleft\arraybackslash}p{0.13\linewidth}}}
    \toprule
    & \multicolumn{4}{c}{\bf Номер сгенерированного массива} & {\bf Среднее} \\
    \cmidrule{2-5}
    {\bf n} & {\bf 1} & {\bf 2} & {\bf 3} & {\bf 4} & {\bf значение} \\
    \midrule
    10 & 45 & 45 & 45 & 45 & 45\\
    100 & 4950 & 4950 & 4950 & 4950 & 4950 \\
    1000 & 499500 & 499500 & 499500 & 499500 & 499500\\
    10000 & 49995000 & 49995000 & 49995000 & 49995000 & 49995000\\
    \bottomrule
\end{tabular}
\caption{Количество сравнений при сортировке методом «пузырька»}
\label{tbl:selection-cmp}
\end{table}

\begin{table}[H]
\centering
\begin{tabular}{*{6}{>{\raggedleft\arraybackslash}p{0.13\linewidth}}}
    \toprule
    & \multicolumn{4}{c}{\bf Номер сгенерированного массива} & {\bf Среднее} \\
    \cmidrule{2-5}
    {\bf n} & {\bf 1} & {\bf 2} & {\bf 3} & {\bf 4} & {\bf значение} \\
    \midrule
    10 & 0 & 45 & 15 & 13 & 18.25\\
    100 & 0 & 4950 & 2391 & 2644 & 2496.25\\
    1000 & 0 & 499500 & 258411 & 248603 & 251628.5\\
    10000 & 0 & 49995000 & 25287653 & 24718524 & 25000294.25\\
    \bottomrule
\end{tabular}
\caption{Количество перемещений при сортировке методом «пузырька»}
\label{tbl:selection-swp}
\end{table}

\begin{table}[H]
\centering
\begin{tabular}{*{6}{>{\raggedleft\arraybackslash}p{0.13\linewidth}}}
    \toprule
    & \multicolumn{4}{c}{\bf Номер сгенерированного массива} & {\bf Среднее} \\
    \cmidrule{2-5}
    {\bf n} & {\bf 1} & {\bf 2} & {\bf 3} & {\bf 4} & {\bf значение} \\
    \midrule
    10 & 31 & 34 & 35 & 34 & 33.5\\
    100 & 606 & 610 & 989 & 886 & 772.75\\
    1000 & 9009 & 9016 & 13229 & 13626 & 11220\\
    10000 & 125439 & 125452 & 183551 & 174421 & 152215.75\\
    \bottomrule
\end{tabular}
\caption{Количество сравнений при быстрой сортировке (рекурсивная реализация)}
\label{tbl:heap-cmp}
\end{table}

\begin{table}[H]
\centering
\begin{tabular}{*{6}{>{\raggedleft\arraybackslash}p{0.13\linewidth}}}
    \toprule
    & \multicolumn{4}{c}{\bf Номер сгенерированного массива} & {\bf Среднее} \\
    \cmidrule{2-5}
    {\bf n} & {\bf 1} & {\bf 2} & {\bf 3} & {\bf 4} & {\bf значение} \\
    \midrule
    10 & 0 & 5 & 3 & 7 & 3.75\\
    100 & 0 & 50 & 141 & 166 & 89.25\\
    1000 & 0 & 500 & 2325 & 2325 & 1287.5\\
    10000 & 0 & 5000 & 30921 & 31108 & 16757.25\\
    \bottomrule
\end{tabular}
\caption{Количество перемещений при быстрой сортировке (рекурсивная реализация)}
\label{tbl:heap-swp}
\end{table}
\vspace{0,2cm}
%\textbf{\vspace{0.2cm} \large Выводы\\}
\paragraph{\vspace{0.2cm} \large Выводы\\}
Количество сравнений в сортировке методом «пузырька» в точности совпадает с теоретическими оценками, количество перемещений действительно пропорционально $\frac{1}{4}(N^2 - N)$.

Метод быстрой сортировки (рекурсивная реализация) тоже отвечает приведённым теоретическим данным.

Нетрудно заметить, что метод быстрой сортировки (рекурсивная реализация) показывает себя лучше на массивах любой длины. На массивах относительно малой длины $(n < 1000)$ разница не так сильно выражена, однако для массивов длины $n \geq 1000$ превосходство быстрой сортировки (рекурсивная реализация) над методом «пузырька» очевидно. Поэтому при сортировке следует отдавать предпочтение именно быстрой сортировке, а не методу «пузырька».

\newpage
\section{Структура программы и спецификация функций}
Ниже приведён полный список функций и их характеристика.
\begin{itemize}
    \item \verb|int cmp(const void *n1, const void *n2)|\\
    Функция-компаратор, указатель на которую подаётся в \verb|qsort|. (сортирует по неубыванию модулей).
    \item \verb|int anticmp(const void *n1, const void *n2)|\\
    Функция-компаратор, возвращающая значения, противоположные \verb|cmp|, указатель на которую подаётся в \verb|qsort|.
    \item \verb|int is_sorted_cmp(const void *n1, const void *n2)|\\
    Функция-компаратор, указатель на которую подаётся в \verb|qsort| функции \verb|is_sorted|. (сортирует по неубыванию).
    \item \verb|void create_random(double *arr, int len)|\\
    Функция, генерирующая по адресу \verb|arr| массив длины \verb|len| из случайных значений. Функция использует библиотечные функции \verb|srand()| и \verb|rand()| для рандомизации и знака. 
    \item \verb|void create_order(double *arr, int len)|\\
    Функция, генерирующая по адресу \verb|arr| массив длины \verb|len|, отсортированный в прямом порядке. Сначала с помощью функции \verb|create_random| по этому адресу генерируется случайный массив, затем сортируется с помощью встроенной функции \verb|qsort| и вспомогательной \verb|cmp|.
    \item \verb|void create_reverse(double *arr, int len)|\\
    Функция, генерирующая по адресу \verb|arr| массив длины \verb|len|, отсортированный в обратном порядке. Сначала с помощью функции \verb|create_random| по этому адресу генерируется случайный массив, затем сортируется с помощью встроенной функции \verb|qsort| и вспомогательной \verb|anticmp|.
    \item \verb|void arrcpy(double *src, double *dest, int len)|\\
    Функция, копирующая значения массива \verb|src| в массив \verb|dest| поэлементно (массивы длины \verb|len|).
    \item \verb|void bubblesort(double *arr, int len)|\\
    Сортирует массив \verb|arr| длины \verb|len| методом «пузырька».
    \item \verb|void sort(double *arr, int left, int right)|\\
    Функция, с помощью которой реализуется быстрая сортировка (рекурсивная реализация).
    \item \verb|void quicksort(ll* arr, int N)|\\
    Сортирует массив \verb|arr| длины \verb|len| быстрой сортировкой (рекурсивная реализация).
    \item \verb|void is_sorted(double *arr_to_check, double *orig_arr, int len)|\\
    Проверяет, является ли массив \verb|arr_to_check| отсортированным в необходимом порядке. Функция поэлементно проверяет соответствие элементов массивов \verb|arr_to_check| и \verb|orig_arr|, а затем поэлементно проверяется равенство модулей элементов в соответствующих позициях.
\end{itemize}
Программа выделяет в динамической памяти массивы для эксперимента, заполняет их тремя вышеописанными способами. После каждой сортировки на стандартный поток вывода печатается количество сравнений и обменов, произведенных в процессе неё. В конце работы вся выделенная память освобождается.
\newpage
\section{Отладка программы, тестирование функций}
Для тестирования сортирующих функций была написана и отлажена специальная тестирующая функция, работающая корректно и при больших значениях $n$. Для массивов малой длины $(n \leq 10)$ корректность сортировок дополнительно проверялась вручную. 

\newpage
\section{Анализ допущенных ошибок}
В ходе выполнения задания были допущены следующие ошибки:
\begin{itemize}
    \item Функция \verb|is_sorted| сначала проводила обычное сравнение соответствующих элементов, но порядок сортировки допускал несколько верных результатов, элементы в которых могли бы различаться знаком. Соответственно, сравнивать необходимо модули чисел.
    \item Функция \verb|is_sorted| не проверяла соответствие элементов одного массива элементам другого. С помощью использования нескольких дополнительных массивов, память под которые объявлялась и освобождалась в теле функции, соответствие элементов двух массивов стало проверяться
    \item Изначально функции \verb|is_sorted| подавались два отсортированных массива, следовательно, исходный массив подвергался изменениям до вызова этой функции, что противоречило постановке задачи. Сортировка оригинального массива была помещена внутрь функции, а исходный массив перестал подвергаться изменениям благодаря использованию нескольких дополнительных массивов, память под которые объявлялась и освобождалась в теле функции
\end{itemize}

\newpage
\begin{raggedright}
\addcontentsline{toc}{section}{Список литературы}
\begin{thebibliography}{2}
\bibitem{first} Кнут Д. Искусство программирования для ЭВМ. Том 3. — М.: Мир, 1978.
\bibitem{second} Вирт Н. Алгоритмы и структуры данных. -- М.: Мир, 2010.
\end{thebibliography}
\end{raggedright}

\end{document}
